\documentclass{article}
\usepackage[margin=1in]{geometry} % Widens the page margins
\usepackage{algorithm}
\usepackage{algpseudocode}
\usepackage{amsmath}
\usepackage{amssymb}

\begin{document}

\section*{System Model}

\paragraph{Failure Model.} 
Crash-recovery; no Byzantine/permanent failures. A crash wipes the node's in-memory \texttt{locks\_map}. Any subtransaction not yet prepared is treated as aborted (matches DB's own recovery behavior --- no reconciliation needed). A prepared subtransaction is recovered by asking the 2PC coordinator for its outcome and applying commit/abort directly --- no predecessor re-check needed. Network is asynchronous: messages may be delayed, lost, duplicated, reordered, not corrupted. Partitions eventually heal.

\paragraph{Execution Model.} 
A subtransaction's handler issues \texttt{get}/\texttt{put} calls against its local store as it runs; each call acquires a lock on its key dynamically, at the point of access, rather than declaring its full read/write set upfront. A subtransaction proceeds speculatively past a conflicting predecessor once that predecessor reaches \texttt{executed} --- its value is fixed even though undecided --- and only blocks on full resolution (\texttt{committed}/\texttt{aborted}) at \texttt{prepare}. A later write to an already-read key triggers a lock upgrade in place, or off the shared node if other readers joined it; if the chain has already moved on past that point, the upgrade cannot be safely spliced in and the subtransaction aborts. Each subtransaction's \texttt{execute} runs its handler code as a single logical unit and ends by flipping status to \texttt{executed} --- but this is not uninterruptible in the scheduling sense: individual \texttt{get}/\texttt{put} calls within it can block on \Call{lock}{} waiting for another subtransaction, and a \texttt{MUST\_ABORT} result from any \texttt{get}/\texttt{put} (whether from \Call{lock}{} or from \Call{upgrade}{}'s case 3) requires the handler to stop executing immediately rather than continue. Application code may not access the store or acquire locks except through \texttt{get}/\texttt{put}.

\paragraph{Architecture.} 
Each participant runs one shim, colocated with its own database. A subtransaction is bound to a single participant's database by definition --- cross-database access is expressed as multiple subtransactions within one Sonata global transaction, coordinated by Sonata's outer 2PC, not by any single subtransaction spanning stores. The shim's \texttt{locks\_map} is therefore scoped to one participant, with a single point of serialization; no cross-node lock coordination is needed within the shim itself. 

\paragraph{Data Model.} 
The shim locks at the granularity of a single record, identified by its primary key $k$; this is the unit \texttt{get}/\texttt{put} and the \texttt{locks\_map} chains operate on. Values are treated as opaque --- the shim does not interpret or type-check them, only serializes access to them via SHARED/EXCLUSIVE locks.

\vspace{1em}

\begin{algorithm}
\caption{Two-Phase Commit Transactional Lock Manager (Part 1: Data Structures \& Execution)}
\small
\begin{algorithmic}[1]
 
\State \textbf{Data Structures:}
\State $\text{lock\_mode} \in \{\text{SHARED}, \text{EXCLUSIVE}\}$
\State $\text{LockNode}: \{\text{mode}, \text{holders}: \text{set of txn}, \text{predecessor}, \text{next}\}$
//TODO: integrate this data structure
\State $\text{locks\_map}\colon \texttt{ConcurrentHashMap<key, LinkedList<LockNode>>}$

\State $\text{locks\_map}: \text{key} \to \text{chain (linked list of LockNode)}$ \Comment{one chain per key, generations of holders}
\State $\text{txn}: \{\text{txn\_id}, \text{locks\_acquired}: \text{key} \to (\text{mode}, \text{target\_node}) \text{ (built dynamically)}, \text{status}, \text{code}\}$
\State $\text{status} \in \{\text{executed}, \text{prepared}, \text{committed}, \text{aborted}\}$
\State \Comment{TODO: intermediate states for retry mechanism}
 
\Statex
\Procedure{execute}{txn}
    \State $txn.\text{code}$ calls \Call{get}{k, txn} or \Call{put}{k, v, txn} multiple times
    \State \Comment{may block inside a get/put on a predecessor; a FAILED result requires the handler to stop calling get/put and return immediately}
    \State \Call{CAS}{\text{txn.status}, \text{NULL}, \text{executed}}
\EndProcedure
 
\Statex
\Procedure{get}{k, txn}
    \State $result \gets$ \Call{lock}{k, \text{SHARED}, txn}
    \If{$result == \text{MUST\_ABORT}$}
        \State \Call{abort}{txn}
        \State \textbf{return NULL, FAILED}
    \EndIf
    \State \textbf{return} value of $k$, SUCCEEDED
\EndProcedure
 
\Statex
\Procedure{put}{k, v, txn}
    \State $result \gets$ \Call{lock}{k, \text{EXCLUSIVE}, txn}
    \If{$result == \text{MUST\_ABORT}$}
        \State \Call{abort}{txn}
        \State \textbf{return FAILED}
    \EndIf
    \State set value of $k$ to $v$
    \State \textbf{return SUCCEEDED}
\EndProcedure
 
\end{algorithmic}
\end{algorithm}
 
\begin{algorithm}
\caption{Two-Phase Commit Transactional Lock Manager (Part 2: Access \& Locking)}
\small
\begin{algorithmic}[1]
 
 
\Statex
\Procedure{lock}{k, mode, txn}
    \If{$\text{txn.status} == \text{aborted}$}
        \State \textbf{return} MUST\_ABORT \Comment{txn was cascade-aborted while running}
    \EndIf
    \State $target\_node \gets \text{null}$
    \If{$\text{txn.locks\_acquired}[k]$ contains $(mode, *)$}
        \State \textbf{return} ALREADY\_HELD
    \ElsIf{$\text{txn.locks\_acquired}[k]$ contains $(\text{EXCLUSIVE}, *)$}
        \State \textbf{return} ALREADY\_HELD
    \ElsIf{$\text{txn.locks\_acquired}[k]$ contains $(\text{SHARED}, *)$ \textbf{and} $mode == \text{EXCLUSIVE}$}
        \State \textbf{return} \Call{upgrade}{k, txn}
    \EndIf
    \State \textbf{atomic}(locks\_map) \{
        \If{$k \notin \text{locks\_map}$}
            \State \indent $\text{locks\_map}[k] = \text{new chain}$
        \EndIf
        \State \indent $chain \gets \text{locks\_map}[k]$
        \State \indent $last\_node \gets chain.\text{last\_node}$
        \If{$last\_node == \text{null}$}
            \State \indent $target\_node \gets \text{LockNode}\{mode, \text{holders}=\{txn\}, \text{predecessor}=\text{null}, \text{next}=\text{null}\}$
            \State \indent append $target\_node$ to $chain$
        \ElsIf{$mode == \text{SHARED}$ \textbf{and} $last\_node.\text{mode} == \text{SHARED}$}
            \State \indent $last\_node.\text{holders.add}(txn)$
            \State \indent $target\_node \gets last\_node$
        \Else
            \State \indent $target\_node \gets \text{LockNode}\{mode, \text{holders}=\{txn\}, \text{predecessor}=last\_node, \text{next}=\text{null}\}$
            \State \indent append $target\_node$ to $chain$
        \EndIf
        \State \indent $\text{txn.locks\_acquired}[k] = (mode, target\_node)$
    \}
    \If{$target\_node.\text{predecessor} \neq \text{null}$}
        \State \textbf{wait until} $\forall t \in target\_node.\text{predecessor.holders}$, $t.\text{status} \in \{\text{committed}, \text{aborted}, \text{executed}, \text{prepared}\}$
        \State \Comment{TODO: no deadlock detection yet --- this wait is unbounded and can deadlock against another txn waiting on this one; accepted for now, to be addressed later}
    \EndIf
    \State \textbf{return} LOCK\_ACQUIRED
\EndProcedure
 
\Statex
\Procedure{upgrade}{k, txn}
    \State $must\_wait \gets \text{false}$
    \State $node \gets \text{null}$
    \State \textbf{atomic}(locks\_map) \{
        \State \indent $node \gets \text{txn.locks\_acquired}[k].\text{target\_node}$
        \If{$node == \text{locks\_map}[k].\text{last\_node}$ \textbf{and} $node.\text{holders} == \{txn\}$}
            \State \indent $node.\text{mode} = \text{EXCLUSIVE}$ \Comment{case 1: pure in-place upgrade}
            \State \indent $\text{locks\_map}[k].\text{last\_node} = node$
            \State \indent \textbf{return} UPGRADED
               \State \indent \Comment{TODO: Concurrent transactions cannot stack TBD}
        \ElsIf{$node == \text{locks\_map}[k].\text{last\_node}$ \textbf{and} $node.\text{holders} \supseteq \{txn\}$}
            \State \indent $node.\text{holders.remove}(txn)$
            \State \indent $new\_node \gets \text{LockNode}\{\text{EXCLUSIVE}, \text{holders}=\{txn\}, \text{predecessor}=node, \text{next}=\text{null}\}$
            \State \indent append $new\_node$ to chain
            \State \indent $\text{txn.locks\_acquired}[k] = (\text{EXCLUSIVE}, new\_node)$
            \State \indent $must\_wait \gets \text{true}$
        \Else
            \State \indent \Comment{TODO: to be discussed whether there is another solution to this}
            \State \indent \textbf{return} MUST\_ABORT \Comment{case 3: something newer already appended}
        \EndIf
    \}
    \If{$must\_wait$}
        \State \textbf{wait until} $\forall t \in node.\text{holders}$, $t.\text{status} \in \{\text{committed}, \text{aborted}, \text{executed}, \text{prepared}\}$
        \State \Comment{TODO: same unbounded-wait / no-deadlock-detection caveat as in \Call{lock}{}}
    \EndIf
    \State \textbf{return} UPGRADED
\EndProcedure
 
\end{algorithmic}
\end{algorithm}
 
\begin{algorithm}
\caption{Two-Phase Commit Transactional Lock Manager (Part 3: Commit Protocol)}
\small
\begin{algorithmic}[1]
 
\Procedure{prepare}{txn}
    \If{$\text{txn.status} \notin \{\text{executed}, \text{prepared}\}$}
        \State \Comment{covers the already-aborted case too: cascade abort sets txn.status directly, so no predecessor dereference is needed here}
        \State \textbf{Send no to 2PC controller}
        \State \textbf{return}
    \EndIf
 
    \For{\textbf{each} $(k, mode, node) \in \text{txn.locks\_acquired}$}
    \If{$node.\text{predecessor} \neq \text{null}$}
        \State \textbf{wait until} $\big(\forall t \in node.\text{predecessor.holders}, t.\text{status} \in \{\text{committed}, \text{aborted}\}\big)$ \textbf{or} $\text{txn.status} == \text{aborted}$
        \State \Comment{txn may have been cascade-aborted by a different key's predecessor while waiting here}
        \If{$\text{txn.status} == \text{aborted}$}
         \State \textbf{Send no to 2PC controller} 
            \State \textbf{return}
        \EndIf
        \If{$\exists t \in node.\text{predecessor.holders}$ with $t.\text{status} == \text{aborted}$}
            \State \Call{abort}{txn}
            \State \textbf{return}
        \EndIf
    \EndIf
\EndFor
 
  \If{not \Call{CAS}{\text{txn.status}, \{executed, prepared\}, prepared}}
  \State \textbf{Send no to 2PC controller} 
  
    return
\EndIf
   \State \textbf{Send yes to 2PC controller}
\EndProcedure
 
\Statex
\Procedure{commit}{txn}
    \If{\textbf{not} \Call{CAS}{\text{txn.status}, \text{prepared}, \text{committed}}}
        \State \textbf{return}
    \EndIf
    \State \textbf{atomic}(locks\_map) \{
        \For{\textbf{each} $(k, mode, target\_node) \in \text{txn.locks\_acquired}$}
            \State \indent remove $txn$ from $target\_node.\text{holders}$
            \If{$target\_node.\text{holders}$ is empty}
                \State \indent\indent remove $target\_node$ from chain
            \EndIf
        \EndFor
    \}
    \State \Comment{safe to evict txn from any external txn-id registry after this point --- nothing in locks\_map references it anymore}
\EndProcedure
 
\Statex
\Procedure{abort}{txn}
    \If{\textbf{not} \Call{CAS}{\text{txn.status}, \{\text{executed},
    \text{prepared},\text{NULL}\}, \text{aborted}}}
        \State \textbf{return}
    \EndIf
 
    \State $to\_cascade \gets \text{set of txn}$
    \State \textbf{atomic}(locks\_map) \{
        \For{\textbf{each} $(k, mode, target\_node) \in \text{txn.locks\_acquired}$}
            \State \indent $entry \gets \text{locks\_map}[k]$
            \If{$entry == \text{null}$}
                \State \indent\indent \textbf{continue}
            \EndIf
            \If{$mode == \text{EXCLUSIVE}$}
                \State \indent\indent $successor\_node \gets target\_node.\text{next}$
                \While{$successor\_node \neq \text{null}$}
                    \State \indent\indent\indent $to\_cascade.\text{update}(successor\_node.\text{holders})$
                    \State \indent\indent\indent $successor\_node \gets successor\_node.\text{next}$
                \EndWhile
            \EndIf
            \State \indent remove $txn$ from $target\_node.\text{holders}$
            \If{$target\_node.\text{holders}$ is empty}
                \State \indent\indent remove $target\_node$ from chain
            \EndIf
        \EndFor
    \}
    \For{\textbf{each} $t \in to\_cascade$}
        \State \Call{abort}{t} \Comment{recursive; CAS makes this idempotent if $t$ already resolved}
    \EndFor
    \State \Comment{safe to evict txn from any external txn-id registry after this point --- status is already set for any successor to observe directly, and a duplicate/late message for an evicted id can be treated as a no-op abort}
\EndProcedure
 
\end{algorithmic}
\end{algorithm}
 
\end{document}
